\section{A relative equivalence in intuitionistic primitive recursive arithmetic}
\label{sec:relative-ipra}

This section proves a relative statement about the delivered machine and
the primary PA194 table.  It does not assume that the primary table has
already been calibrated to PA in iPRA, and it does not prove that missing
calibration.  The proof here is an ordinary mathematical proof with finite
symbolic certificates; it is not a newly compiled Lean or Rocq theorem.

Write $M_0$ for the primary table with SHA-256
\path{4bfea62b17aad7d8c8c5e682fe8942e0a86b93e4351f8e675e32705bdaae007f}
and $M_1$ for PA194AX11QH, whose table has SHA-256
\path{19e215df20429fa6a35f0a99ba0dcfcb40893a6e4b9a4ae6e450b0225f110e50}.
In this section a halting sentence means an existential assertion about
the primitive-recursive iteration of the literal blank-start table,
with the halt state made absorbing.  Thus no external truth predicate
or interpretation of arbitrary PA formulas occurs in the assertion.

\begin{theorem}[Relative machine equivalence]
\label{thm:relative-ipra}
For the two specified tables,
\[
 \iPRA\vdash \Halt(M_0)\leftrightarrow\Halt(M_1),
 \qquad
 \iPRA\vdash \NH(M_0)\leftrightarrow\NH(M_1).
\]
Consequently any iPRA calibration of either table to $\Con(\PA)$
would transfer to the other.  This consequence does not discharge the
calibration premise.
\end{theorem}

The details below explain why the inductions needed for this relative
statement are available in iPRA, instead of merely asserting that all
the algorithms are primitive recursive.

\subsection{The two finite numerical proof predicates}

For clarity, all numerical constructors used here are fixed explicitly.
Let $T(0)=0$, $T(s+1)=T(s)+s+1$, and $p(a,b)=T(a+b)+b$.
Write
\begin{align*}
 I(a,b)&=4p(a,b),& M(a,b)&=4p(a,b)+1,\\
 U(a,b)&=4p(b,a)+2,& E(a,b)&=4p(a,b)+3,\\
 N(a)&=I(a,1),& C(a,b)&=N(I(a,N(b))),\\
 X(a,b)&=N(U(a,N(b))).
\end{align*}
The six heads, in the parameter order used by the emitter, are
\begin{align*}
 L_1(D,H)&=I(X(D,U(D,H)),H),\\
 L_2(D,A,H)&=I(U(D,U(A,H)),U(A,U(D,H))),\\
 L_3(D,A,H)&=I(M(H,D),I(E(H,A),M(A,D))),\\
 L_4(D,A,H)&=I(E(H,D),I(M(A,H),M(A,D))),\\
 L_5(D,H)&=X(D,E(H,D)),\\
 L_6(D,A,H)&=I(U(D,A),I(U(D,I(A,H)),U(D,H))),\\
 L_{6Q}(D,A,H)&=I(U(D,I(A,H)),I(U(D,A),U(D,H))).
\end{align*}
The dominated head and diagonal are
\[
 \rho_8(a,b,c,u)=I(I(I(c,u),a),I(I(a,c),I(b,c))),
 \qquad \Gamma(c)=I(c,U(c,c)).
\]
To define the guarded gate, put
\begin{align*}
 V(a,b,c)&=I(M(a,b),U(c,M(a,b))),\\
 P(a,b,c)&=I(c,M(c,b)),\\
 R(a,b,c)&=U\bigl(b,C(I(M(a,b),P(a,b,c)),I(M(c,b),c))\bigr),\\
 R'(a,b,c)&=U(c,I(R(a,b,c),N(P(a,b,c)))),\\
 G(a,b,c)&=N(I(V(a,b,c),R'(a,b,c))),\\
 \widehat G(a,b,c)&=
   \begin{cases}0,&c=a\text{ or }c=b,\\G(a,b,c),&\text{otherwise.}\end{cases}
\end{align*}
This is the separate-diagonal gate called \texttt{oldBundle} in the
pinned source; it is not the different equality-vacuity bundle also
defined there.  The generated entries are, exactly,
\begin{align*}
 (q_0,q_1,q_2,q_3,q_4)
   &=\bigl(\widehat G(H,D,c),\Gamma(c),L_{6Q}(D,c,H),L_5(D,H),L_4(D,c,H)\bigr),\\
 (q_5,q_6,q_7,q_8,q_9)
   &=\bigl(L_3(D,c,H),L_2(D,c,H),L_1(D,H),\rho_8(D,c,H,d),U(d,h)\bigr),
\end{align*}
with $L_6$ in place of $L_{6Q}$ for the primary machine.

Let $\mathcal C$ be the separate-diagonal numerical calculus of the
primary machine.  It has the six displayed heads, the dominated
implication head $\rho_8$, the guarded gate, the diagonal clause, modus
ponens and theorem generalization.  Let $\mathcal C_Q$ change only
\[
 P\to(Q\to R)\quad\hbox{to}\quad Q\to(P\to R)
\]
in the sixth head.  Here $P=U(D,A)$, $Q=U(D,I(A,H))$, and $R=U(D,H)$.
More precisely, an axiom proof is an instance of $L_1,\ldots,L_6$,
of $\rho_8(a,b,c,u)$ with $u\leq a$, of $G(a,b,c)$ with
$c\ne a,b$, or of $\Gamma(c)$.  The two rules are
$(a,I(a,b))\mapsto b$ and $a\mapsto U(i,a)$ for arbitrary $i$.
This fixes the finite proof checker without appealing to a semantic
interpretation of these numerical expressions.

Proofs can be coded as finite trees with their conclusion included at
each node.  Use a tagged pairing code which is larger than the sum of
its child codes and its data fields.  The assertion that such a tree
is valid, its conclusion, its height, and its number of nodes are all
primitive recursive.  For a valid proof code $D$, every child code,
displayed conclusion, and numerical parameter is at most $D$, and the
height and number of nodes are at most $D$.  These properties are
proved by bounded tree traversal and quantifier-free induction.
Every leaf has height one and an internal node has height one plus
the maximum child height; there are no valid height-zero proofs.

The exchange certificate has 29,932 free implication terms and 18,897
proof rows: 8,986 instances of the \emph{restricted} $\rho_8$ axiom
and 9,911 modus-ponens rows.  The independent check verifies
\[
 (a\to(b\to c))\to(b\to(a\to c)).                 \tag{E}
\]
Each domination condition $u\leq a$ in a $\rho_8$ row is witnessed
by the fact that its term $u$ is a syntactic subterm of its parameter
$a$.  Since $I(x,y)=4p(x,y)$ is at least both $x$ and $y$, each such
inequality has an iPRA proof.  Each MP row is an identity in the free
implication algebra.  Instantiating all 18,897 rows at arbitrary natural
$a,b,c$ therefore gives a fixed iPRA-verifiable proof-producing template.
The certificate does not assume $K$, $S$, unrestricted $\rho_8$, or
quantifier distribution as additional axioms.

Replace a sixth-head leaf by the opposite sixth-head leaf and (E),
followed by MP; keep every other leaf and both proof constructors.
Doing this in either direction is primitive recursive.  Its validity
is verified by bounded induction through the coded input tree: all
intermediate certificates are finite lists, every local validity test
is decidable, and the fixed exchange template contributes no unbounded
proof search.  Thus iPRA proves
\[
 \exists e\,\operatorname{Prf}_{\mathcal C}(e,1)
 \quad\longleftrightarrow\quad
 \exists e\,\operatorname{Prf}_{\mathcal C_Q}(e,1).       \tag{1}
\]
The numerical alias $I(0,0)=0$ is harmless: evaluation of a symbolic
identity preserves it, and the MP predicate tests numerical equality.

\subsection{The exact command interfaces and finite words}

Put $[a;w]=p(a,w)$, with total projections $\operatorname{hd}$ and
$\operatorname{tl}$.  For the new machine the commands, on the store
$(t,w,d,h)$, are
\begin{align*}
 S(t,w,d,h)&=
 \begin{cases}\mathsf{halt},&h=1,\\
 (h,\operatorname{tl}w,d,\operatorname{hd}w),&h\ne1,
 \end{cases}\\
 \operatorname{MP}_x(t,w,d,h)&=
 \begin{cases}(b,w,b,0),&I(t,b)=h,\\(0,w,b,0),&I(t,b)\ne h,
 \end{cases}\qquad b=d+x,\\
 A_{c,x,y,z}(t,w,d,h)&=
       (0,[q_0,q_1,\ldots,q_9,t;w],z,0).
\end{align*}
Here $D=d+x$, $H=h+y$, and the generated entries are the gate,
diagonal, exchanged sixth head, the other five heads, the dominated
head, and $U(d,h)$, in that order.  They are exactly the entries in
the free-constructor emitter certificate.  In particular $q_0,\ldots,q_8$
are zero modulo four and $q_9$ is two modulo four, so none is one.

Let $A_0=A_{0,0,0,0}$.  The following finite words have the indicated
effects for \emph{arbitrary} numerical stacks, including $[0;0]=0$:
\[
 \operatorname{Push}=A_0S^{10}\operatorname{MP}_0,
 \qquad\operatorname{Pop}=A_0S^{12},
 \qquad\operatorname{Discard}=\operatorname{MP}_0.
\]
Push sends $(t,w,d,h)$ to $(0,[t;w],0,0)$.  Discard sends
$(t,w,0,h)$ to $(0,w,0,0)$.  Pop, on $(t,[a;w],d,h)$, halts if
$t=1$ and otherwise returns $(t,w,0,a)$.  All these statements follow
by a fixed number of substitutions in the displayed command equations.
They are quantified equations with decidable case splits and require
no induction over formula truth.

For $0\leq j\leq8$, the selector
\[
 \operatorname{Discard};A_{A,D,H,0};S^{j+2};
                  \operatorname{Pop}^{9-j}                 \tag{2}
\]
returns $(q_j,w,0,0)$ from any $(t,w,0,0)$.  For the dominated
instance $\rho_8(D,A,H,u)$, with $u\leq D$, use
\[
 \operatorname{Discard};\operatorname{MP}_u;
       A_{A,D-u,H,0};S^{10};\operatorname{Pop}.             \tag{3}
\]
The old cache is zero after Discard; consequently the intermediate
MP sets the offset to $u$, including at $u=0$.  Generalization is
\[
 \operatorname{Gen}_i=\operatorname{Push};\operatorname{MP}_i;S;
                A_{i,0,0,0};S^{11},                       \tag{4}
\]
which returns $(U(i,a),w,0,0)$ on $(a,w,0,0)$, even when $a=1$.

On $(0,[a,I(a,b),0;w],0,0)$, the suffix
$S;S;\operatorname{MP}_b;\operatorname{Pop}$ either halts
(when $a=1$ or $b=1$) or returns $(b,w,0,0)$.  Therefore, at an
MP proof node with major proof $e_1$ and minor proof $e_0$, compile
\[
 W_e=\operatorname{Discard};\operatorname{Push};W_{e_1};
       \operatorname{Push};W_{e_0};\operatorname{Push};
                 S;S;\operatorname{MP}_b;\operatorname{Pop}.  \tag{5}
\]
The initial Push reserves the zero needed by the suffix.  Axiom leaves
use (2)--(3), and a generalization node appends (4).  It follows that
\[
 \operatorname{Eval}(W_e,(t,w,0,0))=\mathsf{halt}
 \quad\text{or}\quad
 \operatorname{Eval}(W_e,(t,w,0,0))=(\operatorname{concl}(e),w,0,0).
                                                               \tag{6}
\]
Thus, when $\operatorname{concl}(e)=1$, the word
\[
             V_e=A_0;W_e;\operatorname{Push};S;S            \tag{7}
\]
halts from every carried store.  A lone final Shift would not suffice
for the new cache-only guard.

For completeness, the primary machine uses the same arithmetic heads
before the exchange, but its AX leaves $q_0$ visible and pushes
$q_1,\ldots,q_9,t$; its successful MP additionally pushes a zero;
its Shift tests $t+h=1$.  Its corresponding words are
\begin{align*}
 \operatorname{Push}_0&=A_0S^9\operatorname{MP}_0,&
 \operatorname{Pop}_0&=A_0S^{11},\\
 \operatorname{Discard}_0&=\operatorname{MP}_1;
       A_{1,0,0};S^{10};\operatorname{MP}_0,&
 \operatorname{Offset}_i&=
      \begin{cases}\epsilon,&i=0,\\\operatorname{MP}_i,&i>0.
      \end{cases}
\end{align*}
Here the primary Discard contract is asserted on inputs with
$d=h=0$, exactly as supplied by its compiler; the new Discard's
arbitrary-cache contract is not attributed to the primary word.
Its Pop is total even when $t=1$.  The $q_j$ selector, $1\leq j\leq8$,
uses $\operatorname{Discard}_0;A_{A,D,H};S^{j+1};
\operatorname{Pop}_0^{9-j}$, with (3)'s analogous offset treatment
for a nonzero lower bound.  Its gate selector is
$\operatorname{Discard}_0;A_{c,b,a};\operatorname{Pop}_0^{10}$.
Generalization uses Push, Offset, Shift, $A_{i,0,0}$, and ten Shifts.
At an MP node there is no initial reserved-zero Push: the successful
MP supplies that zero.  The suffix is
$S;S;\operatorname{MP}_b;A_{b,0,0};S^{11}$.
The primary refutation word is $A_0;W_e;S$.
Direct substitutions give the same invariant (6) for its compiler.
These exact identities are also displayed in the repository's
\texttt{PA194Words.lean} and \texttt{PA194DerivationCompiler.lean};
the argument here uses their equations, not an unproved transfer of
an external Lean induction into iPRA.

\subsection{Why the compiler induction is weak enough}

The phrase ``induct on proofs, uniformly in the old stack'' can hide
an induction with unbounded quantifiers.  Here it can be replaced by
an explicitly bounded induction.  Fix a decorated proof-code bound $D$
and an initial bound $B_0\geq\max(D,t,w,1)$.  Define
\[
             B_{j+1}=(4B_j+4)^8.                          \tag{8}
\]
This bounds \emph{inputs of recursive subproof calls}, not every
intermediate value of an AX execution.  Before such a call in (5),
at most two stack pushes have been made, their pushed conclusions
are at most $D$, and each push satisfies
\[
        p(a,w)\leq(a+w+1)^2.
\]
Thus (8) is more than sufficient.  The primary compiler needs no
larger bound.  A successful subproof restores its protected old
stack, so its internal AX values do not enter this bound.

Put $k^*=\min(k,D)$.  Let $P(k)$ say that (6) holds for every valid proof
code at most $D$ of height at most $k^*$, and for all
$t,w\leq B_{D-k^*}$.
All these quantifiers are bounded, and evaluation of a finite word is
primitive recursive.  Thus $P(k)$ has a primitive-recursive characteristic
function and is usable as a quantifier-free induction predicate in
iPRA.  The base case is empty.  At a step from $k$ to $k+1$ with $k<D$, the
recursive calls have height at most $k$ and inputs at most $B_{D-k}$,
by (8).  The fixed leaf and suffix identities prove the remaining
cases.  If a proof has smaller height already, monotonicity of the
bounds applies.  For $k\geq D$ the predicate is unchanged, so those
induction steps are immediate.  At $k=D$ the result is (6) for the originally arbitrary
$t,w$.  This supplies the required uniformity without induction on an
unbounded $\forall w$ formula.

The compiler has at most 256 commands per input tree node, so its
word length is at most $256D$.  This bound is deliberately loose:
the largest leaf selector has fewer than 140 commands, and the
nonrecursive overhead at an MP or generalization node is smaller.
Repeated primitive-recursive concatenation therefore defines the
compiler and its complete evaluator.  No proof search is involved.

\subsection{Soundness by explicit proof annotations}

An annotation consists of a valid numerical proof of $t$, a valid
proof of $h$, and a \emph{tagged finite list} of valid proof codes for
the stack.  Fold the conclusions of that list with $p$ and tail zero
to obtain the numerical stack $w$.  The annotation need not be the
canonical Cantor expansion: a trailing theorem of zero is permitted.
For an empty annotation list, use the fixed proof of zero, which is
the instance $\rho_8(0,0,0,0)=0$.

This choice makes all zero aliases exact.  Pushing a zero proof can
extend the tagged annotation list even when $p(0,0)=0$ leaves the
numerical stack unchanged.  For a nonempty list the Cantor projection
identities give the corresponding head and tail; for an empty list
they give zero and zero.  No fictitious nonzero list length is read
off the number zero.

AX attaches its finite axiom proofs, a generalization of the old
cache proof, and the saved top proof.  A failed guard attaches the
zero proof.  Successful MP combines its two premise proofs; failure
attaches zero.  A primary successful MP also attaches the dummy-zero
proof to the stack.  Shift moves proofs between these positions.
The offset has no proof obligation.  The new halt test reads a proof
of one in the cache.  For the primary test $t+h=1$, a decidable
nonnegative-integer case split finds a proof of one in the top or
cache.

These updates are primitive recursive, including the updating of
the list of proof codes.  Their invariant is the decidable statement
``all listed proof objects pass the proof checker and their folded
conclusions equal the displayed store''.  Quantifier-free induction
on the scheduler clock therefore constructs, from any finite halting
scheduler run, an actual numerical proof of one.  This is a syntactic
proof transformer; it assumes no consistency or soundness of either
numerical calculus.

\subsection{An explicit bound for scheduler fairness}

For a finite word use $E(\epsilon)=0$ and
$E(c::v)=1+b_c(E(v))$.  The new command constructors are
\[
 b_S(z)=p(0,z),\quad
 b_{\operatorname{MP}_x}(z)=p(1,p(x,z)),\quad
 b_{A_{c,x,y,u}}(z)=p(c+2,p(y,p(x,p(u,z)))).                \tag{9}
\]
The primary AX has the same expression with its extra $p(u,z)$
replaced by $z$.  A fresh candidate $n$ is decoded directly; a
positive continuation $l$ is decoded from $l-1$.  Every extracted
continuation is at most its input seed.  Hence a positive continuation
strictly decreases, and candidate $n$ takes at most $n+1$ commands
unless it halts.

Given the all-store refutation word $V_e$, choose the single candidate
\[
 n=p(2,p(0,p(0,p(0,E(V_e)))))\geq3                        \tag{10}
\]
for the new machine; omit the innermost $p(0,\cdot)$ for the primary.
It first executes $A_0$ and then $V_e$.  The prefix may alter the
carried store, since (7) works from every store.  The fresh counter
starts at two and increases by one.  Therefore, by fewer than
\[
                  (n+2)^2                               \tag{11}
\]
command steps, the scheduler has either halted earlier or executed
this halting candidate.  Indeed, the commands in candidates
$2,\ldots,n$ number at most $\sum_{j=2}^n(j+1)<(n+2)^2$.
The comparison and the decreasing-continuation argument are bounded
induction on natural numbers and decidable run predicates.  Thus
(10)--(11) provide a primitive-recursive halting witness, rather
than an unbounded appeal to eventual enumeration.

Combining annotations with the two compilers proves in iPRA
\[
 \Halt(\text{primary scheduler})\leftrightarrow
        \exists e\,\operatorname{Prf}_{\mathcal C}(e,1),
 \qquad
 \Halt(\text{new scheduler})\leftrightarrow
        \exists e\,\operatorname{Prf}_{\mathcal C_Q}(e,1).   \tag{12}
\]

\subsection{Primitive source routines: explicit recurrences}

The source-to-command step is also witnessed by primitive-recursive
traces.  Here are its precise cutpoints and workspace invariant.
At the root the nine registers, in their source order, are
\[
       (R_0,\ldots,R_8)=(n,t,w,0,d,h,0,0,l).
\]
At an AX entry only $R_3=a$ is different.  Write $X=R_6$ and
$Y=R_7$ for the two work registers.  All ordinary arithmetic helper
contracts enter and leave with $X=Y=0$ and preserve every register
not mentioned in their effect.  There is one deliberate exception:
the two halves of the gate have the intermediate workspace
$(R_3,X,Y)=(0,a,a)$, specified below.  Source cutpoints are
\[
\begin{array}{c|rrrrr}
 &\text{root}&\text{AX}&\text{MP}&\text{return jump}&\text{Shift}\\ \hline
 M_0&25&50&3554&3660&3661\\
 M_1&25&50&3572&3659&3660
\end{array}
\]
The return jump has successor 25.  The shared prefix 0--24 contains
the explicit halt at 10 and stack unpairing at 11--24.  Fetch and
dispatch occupy 25--49.  These are source ordinals, not physical
Turing-machine state numbers.

The finite checker \texttt{build\_macro\_plan.py} records all 190 AX
token intervals.  It checks their exact resolved instructions and
all successor labels, that each interval has only its first site
as an external entry, and that its only external exit is the next
token's first site.  Of these intervals, 57 are the literal
instruction incrementing $R_1$; the other 133 match one of nine
pinned primary helper graphs under translation of their internal
labels and exit.  Their contracts are as follows; fields omitted
from a displayed assignment remain fixed.
\[
\begin{array}{c|l}
 \text{source node}&\text{effect from a clean workspace}\\ \hline
 49&(t,w)\mapsto(0,p(t,w))\\
 55&(t,w)\mapsto(d,p(t,w))\\
 58&(t,w)\mapsto(h,p(t,w))\\
 70&(t,w)\mapsto(a,p(t,w))\\
 65&(t,p(v,w))\mapsto(I(v,t),w)\\
 72&(h,l)\mapsto(h+\operatorname{hd}l,\operatorname{tl}l)\\
 75&(t,d,l)\mapsto(I(t,d+\operatorname{hd}l),
                      d+\operatorname{hd}l,\operatorname{tl}l)\\
 142&(a,X,Y)=(a,0,0)\mapsto(0,a,a)
\end{array}
\]
The remaining node 150 has the contract
\[
 (t,d,h,X,Y)=(t,d,h,a,a)\longmapsto(g,0,0,0,0),
 \qquad
 g=\begin{cases}0,&h=a\text{ or }d=a,\\t,&\text{otherwise,}\end{cases}
\]
with $R_3=0$ throughout.  Nodes 55, 58 and 70 are
\emph{push followed by copy}, not mere additions to $t$.
Node 75 reads $d$, pushes the incoming $t$, copies the updated $d$,
and constructs $I(t,d)$; its temporary push is consumed and $w$
is preserved.  The contracts hold for arbitrary numerical $w$,
including zero aliases.  At every use of node 65 the symbolic
emitter certificate supplies a nonempty \emph{tagged} stack prefix;
the equation $p(v,w)$, not a purported numerical list length,
justifies taking its head.  Node 75 supplies its own pushed item.
Their detailed row-cycle derivations are
the pinned \texttt{AXMacroSound}, \texttt{Readers},
\texttt{AXGateCompare} and \texttt{PA194AXGate} modules, with the
following explicit primitive-recursive clocks and induction
interpretation.  Equality of every resolved helper row, including
its exit, is what permits transplantation of those derivations;
an abstract macro name alone would not do so.

The independent \texttt{check\_outside\_ax\_rows.py} checks all
140 retained rows outside AX and every successor under the explicit
relocation.  It removes exactly the old MP dummy-push interval
$[3634,3653)$, redirects success to the retained copy-$d$ block at
3652, and checks the two new Shift blocks against the old clear-$t$
and copy-$h$ blocks.  These 140 rows, nine Shift-prefix rows and
the AX intervals cover all 3,671 new source sites.  Thus the two
finite binding checks together cover the complete source, including
the command dispatcher, not only the constructor expression.

The arithmetic routines have only
draining/copying loops, Cantor pairing and unpairing, multiplication
by four, and destructive comparisons.  These are not unspecified
arithmetic oracles.  For example, the actual pairing loop has length
\[
 P(0)=1,\qquad P(s+1)=P(s)+7s+5.
\]
The actual pairing core on inputs $t,x$ takes $4t+1+P(t+x)$ steps,
and pushing $t$ onto $w$ takes
\[
        3t+7w+4+P(t+w)+3p(t,w).
\]
Writing $a_j=\operatorname{advance}^j(0,0)$ for Cantor's diagonal
enumeration, the unpairing loop takes
\[
 U(n)=1+\sum_{j<n}r(a_j),\qquad
 r(a,b)=\begin{cases}3b+8,&a=0,\\5,&a>0.\end{cases}
\]
Its full restoration adds
$3\operatorname{hd}(n)+3\operatorname{tl}(n)+2$ steps.
The individual copy and times-four loops take $7v+2$ and $6v+1$
steps respectively.  Draining loops take $3v+1$ or $4v+1$.
The clear-only loop takes $2v+1$.  The gate comparison, on
$(t,h,x)$, has the primitive-recursive clock
\begin{align*}
 K_t(0,0)&=2t+3,& K_t(0,x+1)&=3x+5,\\
 K_t(h+1,0)&=2h+6,&K_t(h+1,x+1)&=3+K_t(h,x).
\end{align*}
It clears $h,x$ and preserves $t$ on inequality, but clears $t$
on equality.  This is a recursion through $\min(h,x)$: after
$j$ matching pairs, $h$ and $x$ have each decreased by $j$ and
exactly $3j$ rows have executed.  The two unmatched tails then
use the displayed draining clocks.  Gate node 142 takes $4a+1$
steps; the rest comprises $K_t(h,a)$, a $3a+1$ transfer from
$Y$ to $X$, a $3d+1$ transfer from $d$ to $h$, and the second
comparison.  Its intermediate clean-workspace conditions are
therefore accounted for explicitly.
Each formula comes from concatenating the displayed finite row
cycles.  Fix the input registers and induct on the number of
drained units or diagonal positions, expressing the current
registers and clock by their primitive-recursive iterates.  The
invariant is then quantifier-free.  This is the weak-base version
of the row-cycle arguments, including cases where the reference
Lean proof was written as induction generalizing a second register.
No universal quantifier over arbitrary environments is used as
an induction formula.

Fetch preserves $(t,w,d,h)$.  If $l=0$, its copy-and-restore prefix
takes $7n+4$ rows and leaves $(n,l)=(n+1,n)$, with the second $n$
meaning the original value; if $l>0$, it takes two rows and leaves
$(n,l)=(n,l-1)$.  It then unpairs this seed into $a,l$, with
$X=Y=0$, and uses one or two dispatch rows at 48--49.  This proves
exactly the scheduler seed and operand rules in (9).

The MP prelude reads $x$ into $d$, constructs $I(t,d+x)$ without
changing $w$, and drains this value into $X$.  The unchanged
comparison either reaches the success copy-$d$ block or the return
jump; it clears $t,h,X,Y$ on both paths.  Its equality case copies
$d+x$ to $t$ in $7(d+x)+2$ rows.  The only deletion is the dummy
push on that success path.  Thus its full contract is exactly the
new MP equation above for arbitrary $x,t,w,d,h$.
For its exact comparison clock use
\begin{align*}
 J(0,0)&=2,&J(0,x+1)&=3x+5,\\
 J(h+1,0)&=2h+6,&J(h+1,x+1)&=3+J(h,x).
\end{align*}
Unlike the gate comparison, equality exits immediately instead of
executing a clear-$t$ zero test.  Writing $b=d+x$ and $v=I(t,b)$,
the part after construction takes $3v+1+J(h,v)$ rows, then
$7b+2$ rows on success or zero on failure, and one return jump.

The AX core ends at 3519 with $t=G(H,D,c)$, $a=c$, $d=D$,
$h=H$, clean $X,Y$, and numerical stack
$[q_1,\ldots,q_9,t_{\rm old};w_{\rm old}]$.  Its two gate blocks
end at 3553 with $a=d=h=X=Y=0$ and $t=q_0$.  The last push
ends at 3572 with $t=0$ and the advertised eleven-entry stack.
The following MP block consumes the extra field $z$.  Since its
incoming $t=h=d=0$, it returns $t=0,h=0,d=z$ for both $z=0$
and $z>0$.  This proves the AX11 contract, including arbitrary
values of the fourth field, without assuming that field is zero.

The new Shift prefix at 3660 first clears $t$ in $2t+1$ rows,
then copies $h$ to $t$ in $7h+2$ rows while restoring $h$.
The final two rows set $X=1$ and jump to the unchanged comparison
at zero.  Its input is therefore $(t,h,X)=(h,h,1)$.  It halts
exactly when $h=1$, after $2t+7h+11$ rows in that case.
For $h\ne1$, it reaches the unpairing block at 11 after
$2t+9h+10$ rows, with $t=h_{\rm old}$, $h=X=Y=0$.
Unpairing $w$ then returns to 25 with the new Shift store.
These calculations also show that a command cannot halt inside
an arithmetic helper on the way to its specified endpoint.

Concatenate these traces at the displayed clean cutpoints.  All
segment clocks are fixed finite compositions and sums of the
primitive-recursive recurrences just given, evaluated on the
successive macro stores.  This gives a positive source segment
for every scheduler step and the separate terminating segment
for a halt.  The primary contracts are the same inherited
calculations with its original final AX return, MP dummy push
and sum guard.

For an explicit coarse majorant, an elementary routine with all
its numeric registers bounded by $B\geq1$ has duration and maximum
register bounded by $2^{20}(B+1)^4$, using the above recurrences.
One may iterate this majorant along the fixed finite macro plan,
including the parser and final guard routines.  Its iteration
count is the length of that materialized plan, not one application
per entire AX command.  In particular the nested axiom constructors
are not incorrectly assigned a degree-eight bound.  Summing the
actual recurrence costs is an alternative exact segment clock.

\subsection{Why the literal symbolic certificates are sound in iPRA}

The independent literal checker uses a two-sided finite-word tape
representation with arbitrary zero exterior.  A run atom represents
$0^a$ or $1^a$, where $a$ is an affine expression with nonnegative
coefficients in the gap, register values, and spare-list width.
The list blocks are
\[
                  B(v)=01^{v+1},\qquad C(v)=101^v.
\]
The four generic list traversals are derived from the literal rows
by first checking one arbitrary block, of length $v+2$, with both
external contexts left opaque.  Fix the code of a finite list $L$;
quantifier-free induction on the traversed prefix length concatenates
the recorded block traces, using the primitive-recursive functions
for the corresponding prefix and suffix of this fixed $L$.
This avoids induction on an unbounded assertion about all lists.
The only list-normalization
identity is
\[
                   C(L)1=1B(L),
\]
together with its reversal.  Its empty case is $1=1$, and its cons
case follows from $C(v)1=1B(v)$.  This is again an induction about
prefixes of a fixed primitive-recursively coded word, so is available
in iPRA.

There are precisely three operational symbolic rules.  A literal
row consumes one provably present bit.  A same-state,
same-direction self-loop traverses a homogeneous run of length
$a$; quantifier-free induction on $a$ gives its exact write effect,
duration and displacement.  A checked list traversal concatenates
the block lemma and has the recorded total width as its duration.
All other normalizations are word identities: merging equal-bit
runs, deleting an empty run, reversing a finite list, and absorbing
a finite zero suffix into the blank exterior.  None asserts a
runtime fact without a row or an induction just described.

The finite recorded derivations therefore prove the successor-frame
equation for every source site and both decrement branches, with
arbitrary nine natural registers, arbitrary nonnegative header
padding, and an arbitrary nonempty finite spare list.  The clocks
are the recorded affine expressions; the new table's 4,778 cases
all have positive constant part, at least three.  The frame
successors preserve nonnegative padding and a nonempty spare list;
increments append a spare block.  The fixed blank initialization
is a closed finite calculation of the entire nonzero tape support.
Finite row identities and closed startup calculations are
iPRA-provable without a uniform reflection principle for the checker.

It remains to reflect a physical halt rather than just simulate a
source step.  Augment the nine-register source configuration with
the header-padding counter, signed header/head location, and tagged
spare list used by its physical frame.  Update those auxiliary
fields by the finite
successor-frame contract.  They are initialized by the startup
receipt, and their nonnegativity and nonemptiness are the invariant
just established.  Forgetting the auxiliary fields commutes with
every source step.  Conversely, iterating their prescribed updates
lifts every source run from this initial configuration; no choice
or unbounded search is needed.  This augmentation therefore does
not change halting.  In the following, source configurations include
these fields, so that $F$ is a function of the current configuration,
not of an unspecified execution history.

Distinguish a live source configuration whose next
instruction is \texttt{HALT} from an executed-halt tag.  The former
still has its ordinary live control frame.  Executing that instruction
changes it to a tag containing the entire prehalt configuration;
the tag is thereafter absorbing.  Its physical realization is,
by definition, the endpoint of the recorded halt segment applied to
the preserved prehalt frame.  The finite halt contract proves that
this endpoint has the physical halt state; no prescribed final tape
shape is needed.  The endpoint is primitive recursive because the
segment clock and literal table iteration are primitive recursive.
Thus the halted-tag representation does not identify a prehalt
frame with a physical halt, and its halt correspondence is exact.

Let $F(c)$ be this extended full frame, $\ell(c)\geq1$ the
recorded segment clock for live $c$ (and one for an absorbing tag),
and $c_n$ the primitive source iteration.
Put $L_0=0$ and $L_{n+1}=L_n+\ell(c_n)$ while the source has
not executed its halt instruction, and continue with unit increments
afterwards.  The just-proved frame equations, including the
executed-halt tag equation, give
\[
 L_n\geq n,\qquad
 \operatorname{run}_M(b+L_n)=F(c_n),                       \tag{13}
\]
where $b$ is the fixed initialization time.  A source halt witness
therefore gives a physical halt witness.
Conversely, for a physical halt at time
$t$, bounded search among the first $t+1$ source configurations
either finds a source halt or gives, by (13), a nonhalt frame after
time $t$.  The latter contradicts absorption.  The search is
decidable and bounded, so the argument yields a source witness
intuitionistically; it uses no Markov principle.  Use the same
convention at the command-to-\emph{unaugmented}-source boundary:
a command halt first
executes its source segment and then stores that segment's exact
terminal source configuration in an absorbing tag.  It is not
identified with the live source entry at the halt-test command.
Apply the same positive-segment argument to this boundary, whose
frame is just the displayed nine-register clean command entry.
Then apply the deterministic ghost lift/forget equivalence, and
finally the augmented-source-to-physical argument (13).  In
particular, no physical padding, head position or spare list is
claimed to be recoverable from the scheduler store alone.

This proves each literal machine equivalent in iPRA to its
scheduler.  Equations (1) and (12) now prove the first assertion of
Theorem~\ref{thm:relative-ipra}; intuitionistic contraposition proves
the nonhalting assertion.

\subsection{The remaining absolute calibration}

No part of this proof translates an arbitrary internal PA proof
into a numerical-calculus refutation, or a numerical refutation
into a PA refutation.  The pinned repository does establish the
one-sided result $\Con(\PA)\to\NH(\text{modeled scheduler})$;
its README distinguishes that result from a physical-machine
calibration.  The missing direction needed for the full calibration
is, in particular, a verified internal transformation from arbitrary
PA refutations into numerical-calculus refutations.  The ordinary
PA calibration and the conditional iPRA endpoint do not supply
that transformation with its weak-base correctness proof.  Thus
the strengthened full equivalence requested by the user remains
open even though the relative theorem and the new literal
194-state construction are established here.
